\section{Propagation from a faithful component extension}
\label{sec:propagation}

This section explains how a cycle in an extension of a simple component
produces nonzero homology for the ambient group.  The local cycle need not
lie in the largest dimension of the extension.  What matters is a nonzero
coefficient on a flag that contains every rank below its final subgroup.
We first establish the group-theoretic and chain-level facts that make
this coefficient survive passage to the ambient group.

The use of shuffles and full chains to propagate homology is developed
in \citet[Section~3, Lemmas~3.12--3.14]{Piterman2021}, following
\citet[Lemmas~0.24, 0.25, and~0.27]{AschbacherSmith1993}.
Maximal faithful extensions and deletion of faithful overgroups also
occur in \citet[Section~4, proof of Theorem~1, and Section~5]{Piterman2021}.
We give the version needed here, with the supported coefficient,
centralizer condition, and permitted elementary enlargements explicit.
The full-chain method already permits propagation without the Quillen
dimension property; the construction below retains that feature.

\subsection{Centralizers and faithful configurations}

We use two standard consequences of the component theorem: distinct
components commute, and every component centralizes the Fitting subgroup
\cite{GLS1998}.  In particular, distinct conjugates of a nonabelian simple
component have trivial intersection.  The following consequence does not
require the component to be normal.

\begin{lemma}\label{lem:component-core}
Let $G$ be a finite group, let $p$ be a prime, and let $L$ be a nonabelian
simple component of $G$.  If $\Op(G)=1$, then
\[
  \Op(C_G(L))=1.
\]
\end{lemma}

\begin{proof}
List the distinct $G$-conjugates of $L$ as $L=L_1,L_2,\ldots,L_t$.
The component theorem gives an internal direct product
$N=L_1\cdots L_t$, and conjugation by $G$ permutes its factors.  Thus
$N\trianglelefteq G$.

Set $C=C_G(L)$ and $Q=\Op(C)$.  For $i>1$, the group $L_i$ is contained
in $C$.  It is also subnormal in $C$: intersect a subnormal chain from
$L_i$ to $G$ with $C$.  Consequently $L_i$ is a component of $C$.
The normal $p$-subgroup $Q$ is nilpotent, so it lies in the Fitting
subgroup of $C$, and the component theorem inside $C$ gives
$[Q,L_i]=1$.  By definition $Q$ also centralizes $L_1$.  Therefore
\[
 Q\le D:=C_G(N)\le C.
\]
Since $Q\trianglelefteq C$, the inclusion $Q\le D$ implies
$Q\trianglelefteq D$.  Hence $Q\le\Op(D)$.  Finally,
$D\trianglelefteq G$ and $\Op(D)$ is characteristic in $D$, so
$\Op(D)\le\Op(G)=1$.
\end{proof}

\begin{definition}\label{def:faithful-configuration}
Fix a finite group $G$, a prime $p$, and a nonabelian simple component
$L$ of $G$.  A \emph{faithful configuration for $L$ in $G$} is a subgroup
$A\in\Ap(G)$ such that
\[
 A\le N_G(L),\qquad A\cap L\ne1,\qquad
 C_{LA}(L)=1,\qquad \Op(C_G(LA))=1.
\]
We write $H=LA$.  The condition $C_H(L)=1$ says that the conjugation
homomorphism $H\longrightarrow\Aut(L)$ is injective.  We may therefore
identify $L\le H\le\Aut(L)$ using this homomorphism.
\end{definition}

If $\Op(G)=1$, every nonidentity elementary abelian subgroup $A\le L$
gives a faithful configuration: here $LA=L$, the simplicity of $L$
gives $Z(L)=1$, and Lemma~\ref{lem:component-core} gives the last
condition.  This observation will provide initial members of the
families used later.

The next observation separates faithfulness of an elementary subgroup
from faithfulness of the extension it generates with $L$.

\begin{lemma}\label{lem:faithfulness-upgrade}
Let $A$ be a faithful configuration for a nonabelian simple component
$L$ of $G$.  Suppose every element of order $p$ in $C_L(A)$ belongs
to $A$.  If $D\in\Ap(G)$ contains $A$, then $D\le N_G(L)$.
Moreover,
\[
 C_D(L)=1\quad\Longrightarrow\quad C_{LD}(L)=1.
\]
\end{lemma}

\begin{proof}
Choose $1\ne x\in A\cap L$.  For $d\in D$, commutativity of $D$
gives $x^d=x$, so $x\in L\cap L^d$.  Distinct conjugates of the simple
component $L$ have trivial intersection.  Therefore $L^d=L$, as required.

Assume $C_D(L)=1$.  It suffices to prove that every $d\in D$ inducing
an inner automorphism of $L$ already belongs to $L$.  Such an element
can be written
\[
 d=lc,\qquad l\in L,\quad c\in C_G(L).
\]
The factors commute and $L\cap C_G(L)=Z(L)=1$, so this factorization
is unique.  Since $d^p=1$, we have $l^p=c^p=1$.  Every $a\in A$
normalizes both $L$ and $C_G(L)$; comparing the two factorizations of
$d^a=d$ gives $l^a=l$ and $c^a=c$.  Thus $l\in C_L(A)$.
Either $l=1$, or $l$ has order $p$, in which case the hypothesis gives
$l\in A$.  In either case $l\in D$, and consequently
$c=l^{-1}d\in C_D(L)=1$.  Hence $d=l\in L$.

Now an element $ld\in C_{LD}(L)$ makes the action of $d$ inner.
The preceding argument gives $d\in L$, and then
$ld\in L\cap C_G(L)=1$.  This proves the asserted faithfulness of $LD$.
\end{proof}

\subsection{Deletion and augmented chains}

For a finite poset $X$, write $X_{>x}$ and $X_{<x}$ for its strict
upper and lower intervals.  All chains below use rational coefficients
and the augmented simplicial boundary.  In particular, the empty chain
$[\,]$ has degree $-1$, and the boundary of a vertex is $[\,]$.

We shall use the following elementary deletion facts.  If $X_{>x}$ is
nonempty and contractible, then
\[
 \Delta(X\setminus\{x\})\longrightarrow\Delta(X)
\]
is a homotopy equivalence.  Indeed, the link of $x$ is
$\Delta(X_{<x})*\Delta(X_{>x})$, which is contractible, also when
$X_{<x}$ is empty.  The complex $\Delta(X)$ is obtained by attaching
the cone on this link to $\Delta(X\setminus\{x\})$ along the link.
Attaching a cone along a contractible subcomplex preserves homotopy
type, by the homotopy extension property for simplicial complexes.

Also, if a finite group $T$ has $\Op(T)\ne1$, then $\Ap(T)$ is
contractible.  Here is a verification that includes the elementary
subgroup condition.  Let $\mathcal S_p(T)$ be the poset of all
nonidentity $p$-subgroups of $T$.  Writing $Q=\Op(T)$, the order maps
\[
 P\longmapsto P,\qquad P\longmapsto PQ,\qquad P\longmapsto Q
\]
on $\mathcal S_p(T)$ satisfy $P\le PQ\ge Q$, so
$\mathcal S_p(T)$ is contractible.  For any $R\in\mathcal S_p(T)$,
the inverse image of $(\mathcal S_p(T))_{\le R}$ under the inclusion
$\Ap(T)\hookrightarrow\mathcal S_p(T)$ is $\Ap(R)$.
Choose a central subgroup $Z\le R$ of order $p$.  For $E\in\Ap(R)$,
the product $EZ$ is elementary abelian, and the order maps
$E\le EZ\ge Z$ contract $\Ap(R)$.  The finite-poset fiber theorem
\cite{Quillen1978}, in the form asserting that contractible inverse
images of all lower principal ideals give a homotopy equivalence,
therefore applies to this inclusion.  This proves the claim.

The chain operation that will detect the propagated cycle is a residue
at an initial flag.  Its compatibility with boundaries depends on
including every possible rank in that flag.

\begin{lemma}[Saturated residue]\label{lem:saturated-residue}
Let $X$ be a subposet of $\Ap(G)$, and suppose $X$ contains the flag
\[
 \sigma=[A_1<\cdots<A_r=A],\qquad \rk(A_i)=i.
\]
Define a linear map
\[
 R_\sigma:\widetilde C_n(\Delta X;\Q)
       \longrightarrow
       \widetilde C_{n-r}(\Delta X_{>A};\Q)
\]
as follows: a chain containing $\sigma$ is sent to its strict
continuation after $A$, and every other chain is sent to zero.
The continuation of $\sigma$ itself is $[\,]$.  Chain groups in degrees
below $-1$ are zero.  Then
\begin{equation}\label{eq:residue-boundary}
 R_\sigma\partial=(-1)^r\partial R_\sigma.
\end{equation}
In particular, residues of cycles are cycles, and residues of boundaries
are boundaries, including in degree $-1$.
\end{lemma}

\begin{proof}
There is no vertex strictly below $A_1$, and no vertex between
$A_i$ and $A_{i+1}$: these would be elementary abelian subgroups of
ranks strictly between consecutive integers.  Thus every chain
containing $\sigma$ begins with exactly these $r$ vertices.
For such a chain, deleting one of its first $r$ vertices removes part
of $\sigma$ and contributes zero to the residue.  Deleting its
$(r+j+1)$-st vertex, with $j\ge0$, is precisely deletion of the
$(j+1)$-st vertex of its continuation, with boundary sign
$(-1)^{r+j}$ in place of $(-1)^j$.  This proves
\eqref{eq:residue-boundary} on chains containing $\sigma$.
A face of a chain that does not contain $\sigma$ cannot contain
$\sigma$, proving the identity on all remaining chains.  If the
continuation is a single vertex, its boundary is $[\,]$; if it is
empty, its boundary is zero.  The same calculation therefore proves
the identity in the augmented degrees as well.
\end{proof}

\subsection{Products of cycles}

Suppose $H,K\le G$ commute and $H\cap K=1$.  Set
$X=\Ap(H)$, $Y=\Ap(K)$, and adjoin an identity subgroup to each,
writing $X^+=X\cup\{1\}$ and $Y^+=Y\cup\{1\}$.  The poset
\[
 \mathcal P(H,K)
   =\{UV:U\in X^+,\ V\in Y^+,\ (U,V)\ne(1,1)\}
\]
is identified with $(X^+\times Y^+)\setminus\{(1,1)\}$.
Indeed, $UV\cap H=U$ and $UV\cap K=V$, so both the subgroup and
its inclusion relations recover the two coordinates.

\begin{lemma}[Product shuffle]\label{lem:product-shuffle}
With this notation there is a bilinear operation
\[
 \boxtimes:\widetilde C_a(\Delta X;\Q)
       \otimes\widetilde C_b(\Delta Y;\Q)
       \longrightarrow
       \widetilde C_{a+b+1}(\Delta\mathcal P(H,K);\Q)
\]
for $a,b\ge-1$, satisfying
\begin{equation}\label{eq:product-boundary}
 \partial(u\boxtimes v)
   =(\partial u)\boxtimes v+(-1)^{a+1}u\boxtimes\partial v.
\end{equation}
It is the chain construction for the subdivision of the join
$\Delta X*\Delta Y$ by the product poset.  The empty chain is an
identity for this operation.
\end{lemma}

\begin{proof}
Let
$u=[U_1<\cdots<U_r]$ and $v=[V_1<\cdots<V_t]$, where
$r=a+1$ and $t=b+1$, and set $U_0=V_0=1$.
A shuffle is a word with $r$ letters $H$ and $t$ letters $K$.
Starting at $(0,0)$, read the word from left to right, increasing the
first coordinate at an $H$ and the second at a $K$.
At each of the $r+t$ steps record the subgroup $U_iV_j$.
The recorded subgroups form a strict chain.  Give this chain the sign
$(-1)^{\nu}$, where $\nu$ is the number of pairs in which a $K$
precedes an $H$, and define $u\boxtimes v$ to be the sum over all
shuffles with these signs.  When either flag is empty this is just
the other flag, in its pure coordinate; when both are empty it is
$[\,]$.

For the boundary formula, an omitted intermediate vertex between
steps of different types has a second occurrence obtained by
interchanging those two steps.  Their shuffle signs are opposite,
so these occurrences cancel.  This includes deletion of the first
vertex when the first two steps have different types.  The remaining
faces are precisely shuffles in which a vertex of one input flag is
omitted: either two consecutive steps of that type have been joined,
or the final step has been removed.  If the $i$-th $H$-vertex is
preceded by $b'$ $K$-steps, its face sign is $(-1)^{i+b'-1}$,
and removing its $H$-step decreases the inversion count by $b'$.
Its total sign is therefore $(-1)^{i-1}$ relative to the shortened
shuffle.  If the $j$-th $K$-vertex is preceded by $a'$ $H$-steps,
its face sign is $(-1)^{a'+j-1}$, and removing its $K$-step
decreases the inversion count by $r-a'$.  Its total sign is
$(-1)^{r+j-1}$.  Summing these faces gives
\eqref{eq:product-boundary}, since $r=a+1$.  The cases of empty or
single-vertex inputs use the same rule and the augmented boundary.

For completeness, the underlying join description follows from the
usual staircase triangulation of a product of simplices, applied
compatibly to the two cones $\Delta X^+$ and $\Delta Y^+$.
The resulting triangulation is
$\Delta(X^+\times Y^+)$.  Let $s,t\in[0,1]$ be the radial
coordinates in these cones, with zero at the adjoined identities.
Under the product realization, the subcomplex obtained by deleting
$(1,1)$ is exactly the locus $\max(s,t)=1$: in a chain missing the
least element, at least one coordinate is nonidentity throughout.
Conversely, a point with $\max(s,t)=1$ has zero barycentric
coefficient at $(1,1)$ in every containing product simplex.
Radial rescaling replaces this locus by $s+t=1$, the realization of
$\Delta X*\Delta Y$.  This also gives the stated description when
one factor is empty; when both are empty, both complexes are empty.
\end{proof}

\subsection{The propagation lemma}

We can now formulate the precise local input needed from the later
classical-group constructions.  The two maximality conditions in the
statement will ensure that all faithful elementary overgroups can be
deleted, while products with the centralizer remain.

\begin{lemma}[Propagation]\label{lem:propagation}
Let $G$ be a finite group and $p$ a prime.  Assume that every proper
subgroup $T<G$ with $\Op(T)=1$ satisfies
\[
 \widetilde H_*(\Delta\Ap(T);\Q)\ne0.
\]
Let $L$ be a nonabelian simple component of $G$ with $p\mid |L|$,
and let $A$ be a faithful configuration for $L$ in $G$.
Write $H=LA$ and $r=\rk(A)$.  Assume the following three conditions.
\begin{enumerate}[label=\textup{(\roman*)}]
\item Every element of order $p$ in $C_L(A)$ belongs to $A$.
\item There is no $D\in\Ap(G)$ with $D>A$ for which $LD$ acts
faithfully on $L$ and $\Op(C_G(LD))=1$.
\item There is a cycle
$\alpha\in\widetilde C_{r-1}(\Delta\Ap(H);\Q)$ whose coefficient
on some flag
$\sigma=[A_1<\cdots<A_r=A]$, with $\rk(A_i)=i$, is nonzero.
\end{enumerate}
Then $\widetilde H_*(\Delta\Ap(G);\Q)\ne0$.
\end{lemma}

\begin{proof}
We first delete certain vertices without changing the homotopy type,
then construct a cycle in the remaining poset and detect it by
Lemma~\ref{lem:saturated-residue}.

By Lemma~\ref{lem:faithfulness-upgrade}, every elementary overgroup
of $A$ normalizes $L$, and faithfulness of such an overgroup on $L$
implies faithfulness of its extension by $L$.  Define
\[
 \mathcal D=\{D\in\Ap(G):D>A,\ C_D(L)=1\}.
\]
For $D\in\mathcal D$, condition~\textup{(ii)} therefore gives
\begin{equation}\label{eq:deleted-core}
 \Op(C_G(LD))\ne1.
\end{equation}
Delete the elements of $\mathcal D$ in an order in which every strict
overgroup is deleted before its subgroups.

Consider the moment at which $D\in\mathcal D$ is deleted, and
let $U_D$ be its upper interval in the poset then remaining.
Every $E\in U_D$ has $C_E(L)\ne1$, since otherwise $E$ would have
already been deleted.  Put $K_D=C_G(LD)$.  There are order maps
\[
 q_D:U_D\longrightarrow\Ap(K_D),\quad E\longmapsto C_E(L),
 \qquad
 j_D:\Ap(K_D)\longrightarrow U_D,\quad P\longmapsto DP.
\]
The first map is well-defined because $E$ is abelian and contains
$D$, so $C_E(L)$ centralizes both $L$ and $D$.  For the second,
$P$ centralizes $D$ and
$P\cap D\le C_D(L)=1$; hence $DP$ is elementary abelian and
strictly contains $D$.  It centralizes $L$ on the nontrivial subgroup
$P$, so it has not been deleted.  Finally,
\[
 q_Dj_D(P)=P,\qquad j_Dq_D(E)=D C_E(L)\le E.
\]
For the equality, an element $dp\in DP$ centralizes $L$ exactly
when $d$ does, and $C_D(L)=1$.  The displayed comparisons give a
poset homotopy equivalence between $U_D$ and $\Ap(K_D)$.
By \eqref{eq:deleted-core} the latter is contractible.  Each deletion
therefore preserves homotopy type.  Denote the final poset by
\[
 X=\Ap(G)\setminus\mathcal D.
\]

Set $K=C_G(H)$.  It is a proper subgroup of $G$: otherwise it would
contain $L$, contradicting the noncommutativity of $L\le H$.
The configuration gives $\Op(K)=1$, so the inductive assumption
provides a cycle
\[
 \beta\in\widetilde C_b(\Delta\Ap(K);\Q),\qquad b\ge-1,
 \qquad [\beta]\ne0.
\]
Moreover $H\cap K=1$, since an element of this intersection
centralizes $L$ and lies in the faithfully acting group $H$.
Thus Lemma~\ref{lem:product-shuffle} gives a cycle
\[
 z=\alpha\boxtimes\beta
     \in\widetilde C_{r+b}(\Delta\Ap(G);\Q).
\]
In fact $z$ is supported in $X$.  A product vertex $UV$ with
$1\ne V\le K$ has $C_{UV}(L)\ne1$ and so is not deleted.
A pure vertex $U\le H$ could be deleted only if $U>A$.
But such a $U$ would satisfy $LU=H$, making $LU$ faithful with
$\Op(C_G(LU))=\Op(K)=1$, contrary to condition~\textup{(ii)}.
This proves the assertion about the support of $z$.

All elements $E$ of $X_{>A}$ have $C_E(L)\ne1$.  The same maps as
above, now with $D=A$, give
\begin{equation}\label{eq:upper-retraction}
 q:X_{>A}\longrightarrow\Ap(K),\quad E\longmapsto C_E(L),
 \qquad
 j:\Ap(K)\longrightarrow X_{>A},\quad P\longmapsto AP.
\end{equation}
Indeed, $C_E(L)$ also centralizes $A$, while $P\cap A=1$ by
faithfulness; the identities are $qj=\id$ and $jq(E)\le E$.

Let $\lambda\ne0$ be the coefficient of $\sigma$ in $\alpha$.
We claim that the residue of $z$ is
\begin{equation}\label{eq:propagation-residue}
 R_\sigma z=\lambda j_*\beta.
\end{equation}
Every vertex of $\sigma$ lies in $H$.  Because the product
coordinates are unique, a shuffled flag containing $\sigma$ has
trivial $K$-coordinate until it reaches $A$.  It must therefore
take all $r$ of its $H$-steps first, and its $H$-flag must be
exactly $\sigma$: the homogeneous chain $\alpha$ has precisely
$r$ vertices in each term.  The remaining vertices are
$AV_1<\cdots<AV_{b+1}$ for a term of $\beta$.
The shuffle taking every $H$-step first has sign $+1$, which proves
\eqref{eq:propagation-residue}.

If $z=\partial w$ in $\widetilde C_*(\Delta X;\Q)$, then
Lemma~\ref{lem:saturated-residue} makes $R_\sigma z$ a boundary
in $\Delta X_{>A}$.  Apply the augmented chain map induced by $q$,
where a chain with repeated image vertices is sent to zero.
Since $qj=\id$, equation~\eqref{eq:propagation-residue} implies
that $\lambda\beta$ is a boundary in $\Delta\Ap(K)$, contrary
to $[\beta]\ne0$.  Thus $[z]\ne0$.  The homotopy equivalence
$\Delta X\simeq\Delta\Ap(G)$ proves the conclusion.

To make the empty case explicit, $b=-1$ with $[\beta]\ne0$ is
possible only when $\Ap(K)$ is empty.  The map $q$ in
\eqref{eq:upper-retraction} then forces $X_{>A}$ to be empty as
well.  Equation~\eqref{eq:propagation-residue} is a nonzero
multiple of $[\,]$ in that empty upper interval, and cannot be
a boundary.  Hence the same proof applies without requiring a
nonidentity elementary abelian subgroup in $K$.
\end{proof}

\subsection{Choosing a configuration of maximal rank}

In the applications, the cycle may be supported at a different
configuration from the one first selected.  The following criterion
records exactly what must be retained when making that replacement.

\begin{corollary}\label{cor:maximal-family}
Assume the inductive hypothesis on proper subgroups of $G$ in
Lemma~\ref{lem:propagation}, and let $L$ be a nonabelian simple
component with $p\mid|L|$.  Let $\mathcal F$ be a nonempty family
of faithful configurations for $L$ in $G$ such that
\[
 A\in\mathcal F,\quad D\in\Ap(G),\quad D>A,
 \quad D\text{ a faithful configuration}
 \quad\Longrightarrow\quad D\in\mathcal F.
\]
Let $r$ be the largest rank of a member of $\mathcal F$.
If some $A\in\mathcal F$ of rank $r$ is the final subgroup of a
saturated flag having nonzero coefficient in a cycle of degree $r-1$
in $\Delta\Ap(LA)$, then
$\widetilde H_*(\Delta\Ap(G);\Q)\ne0$.
\end{corollary}

\begin{proof}
If $x\in C_L(A)$ has order $p$ and $x\notin A$, then
$D=\langle A,x\rangle$ is elementary abelian and strictly contains
$A$.  It normalizes $L$, has nontrivial intersection with $L$, and
satisfies $LD=LA$.  Thus it is a faithful configuration, so belongs
to $\mathcal F$ by the stated closure property.  Its larger rank
contradicts the choice of $r$.  This proves condition~\textup{(i)}
of Lemma~\ref{lem:propagation}.  Every strict elementary overgroup
of $A$ normalizes $L$, by the first part of
Lemma~\ref{lem:faithfulness-upgrade}; if it also satisfies the two
conditions in~\textup{(ii)}, it is a faithful configuration and
again belongs to $\mathcal F$, contradicting maximal rank.
Thus~\textup{(ii)} holds, and the coefficient hypothesis is
exactly~\textup{(iii)}.
\end{proof}

For example, fix a nonempty collection of nonidentity elementary
abelian subgroups of $L$, and take the configurations containing at
least one of them.  This family has the required closure property,
and Lemma~\ref{lem:component-core} proves its nonemptiness when
$\Op(G)=1$.  If further restrictions are imposed on a family, its
closure under qualifying overgroups must be verified as well.
After constructing a cycle in $H=LA$, it is sufficient to find a
supported subgroup $A'\in\mathcal F$ with $\rk(A')=r$ and
$LA'=H$.  These conditions ensure that the cycle lies in the
required extension and that the maximality argument applies to its
actual nonzero coefficient.  No assertion that $r$ is the full
$p$-rank of $H$ or of $G$ is needed.
